\chapter{Analysis and Discussion}
\label{ch:disc}

\section{What the Results Establish}
\label{sec:interpretation}

The retained eight-view HO-Cap evaluation covers 28 selected
pick-and-place sequences from nine subjects. Prepared 21-joint hand
predictions match all 623,826 valid reference joint slots, with
10.79~mm matched-joint-weighted MPJPE against the dataset's published
semi-automatic MANO fits. This is an external-system perception
measurement, not evidence that the local Azure Kinect rig has the
same accuracy. The object evaluation covers 112 named tracks, but it
uses published instance masks and a reference-assisted first-frame
SE(3) alignment. Its 5.44~mm translation and 7.58-degree rotation
macro-means describe relative tracking drift conditional on those
inputs, not absolute object localization or autonomous segmentation.

\section{Error Sources and Reference Frames}
\label{sec:errbudget}

The hand and object numbers cannot be added into an end-to-end geometric
error budget. The hand score is an absolute comparison in the HO-Cap
world frame; the object score removes a per-object initial gauge
difference. Both refer to semi-automatic annotations rather than an
independent metrology system. The object tracker also receives masks
generated outside ViKi, and static intervals are included alongside
manipulation. Thus a small average object drift need not imply a
sufficiently accurate grasp pose during contact.

The outliers make this distinction operational. The projector remote
\texttt{G18\_2} has 112.60-degree mean rotation error in one sequence
at 98.63\% coverage, but only 6.97 degrees in another. The dandelion
tea and toothpaste box have additional large orientation errors.
The sauce bottle combines 2.82~mm mean translation error with
56.57-degree mean orientation error. These tracks remain in the
aggregate rather than being removed as failures. Symmetry-aware
orientation scoring was not used; the visual review of the remote
did not uniquely diagnose the failure. Conversely, the electric
screwdriver is tracked for only 48.80\% of valid reference frames,
so its conditional error does not describe the missing half.

For the local two-Kinect system, camera geometry, RGB--depth alignment,
exposure timing and table-height registration remain unquantified in
this manuscript. The former local numerical summaries cannot be
rechecked because their source recordings and derived artifacts are
unavailable. Code paths and visually plausible clouds are not
substitutes for a preserved evaluation protocol.

\section{Statistical Validity and Attribution}
\label{sec:validity}

The HO-Cap selection is a single-hand pick-and-place subset, not the
official held-out benchmark split. Subjects and hands are unevenly
represented: 24 right-hand and four left-hand sequences. Joint slots
within the same frames and object frames within the same sequence
are correlated, so their large counts do not imply that many
independent trials. The two-sequence camera-pair pilot was selected
on the same development frames and cannot be interpreted as held-out
evidence. The main eight-view aggregate is descriptive for this
selected cohort; no causal claim about the benefit of camera count,
training data or robot embodiment follows from it.

No surviving local cube dataset supports the previous downstream
figures. In particular, there is no verifiable local stage yield,
simulated stacking rate, policy comparison or contact-safety result.
The current manuscript therefore does not claim an end-to-end
camera-to-simulation proof of concept or Real-to-Sim-to-Real transfer.

\section{Scope and Implications}
\label{sec:implications}

The supported contribution is a reproducible external perception
assessment, with explicit object identities, coverage, outliers and
metric limitations, alongside an implemented offline software
pipeline. The intended cube-transfer experiment remains a research
question rather than a completed result. A new preserved local
cohort should include measurement-only controls: known geometry at
table and carry heights, a marked asymmetric object for orientation,
and disjoint human trials. Retained simulator rollouts and matched
no-demonstration baselines would then be needed before reporting
physical task success or learning benefit. No physical UR10
evaluation is part of the current evidence.
